\documentclass[tikz,border=4pt]{standalone}
\usepackage{ctex}
\usepackage{amsmath}
\usepackage{pgfplots}
\pgfplotsset{compat=1.18}
\usetikzlibrary{calc, arrows.meta}

\begin{document}
% part14/05 新定义题型：图示示意（以某对称函数为例）
% 示意：一次函数和二次函数综合图，展示定义域与值域
\begin{tikzpicture}
\begin{axis}[
  axis lines=center,
  xlabel={$x$}, ylabel={$y$},
  every axis x label/.style={at={(axis cs:5,0)}, anchor=west},
  every axis y label/.style={at={(axis cs:0,7)}, anchor=south},
  xmin=-3, xmax=5, ymin=-3, ymax=7,
  xtick={-3,-2,-1,0,1,2,3,4},
  ytick={-2,-1,0,1,2,3,4,5,6},
  tick label style={font=\small},
  width=7.5cm, height=7.5cm,
  thick,
]
  % 分段函数：x<0时 y=-x，x>=0时 y=x^2-2x+1=(x-1)^2
  \addplot[blue, thick, domain=-2.5:0, samples=30] {-x};
  \addplot[blue, thick, domain=0:3.5, samples=60] {(x-1)^2};
  \node[blue, right, font=\small] at (axis cs:2.8,4.2) {$y=(x-1)^2$};
  \node[blue, above right, font=\small] at (axis cs:-2.5,2.5) {$y=-x$};

  % 分界点
  \fill (axis cs:0,0) circle(2.5pt);
  \node[below right, font=\small] at (axis cs:0,0) {$O$};

  % 顶点
  \fill (axis cs:1,0) circle(2.5pt);
  \node[below, font=\small] at (axis cs:1,0) {$(1,0)$};

  % 标注
  \node[font=\small, gray] at (axis cs:-1.5,-1.8) {$x<0$ 段};
  \node[font=\small, gray] at (axis cs:3.0,1.0) {$x\geq0$ 段};
\end{axis}
\end{tikzpicture}
\end{document}
